\chapter{Preliminaries}
\label{ch:prelim}

This chapter introduces the concepts and tools used in the following chapters.
It starts from the speech signal and acoustic features, moves to the components
of the cascade and then to the architecture of end-to-end speech-to-speech
models, and ends with the measures and statistical concepts of evaluation.

\section{The Speech Signal and Acoustic Features}
\label{sec:features}

In this project speech is represented as a mono signal sampled at 16~kHz with
16-bit integer samples. The continuous signal $x[n]$ is divided into fixed-length
\emph{frames}; in this project a frame is 20 ms long, equivalent to 320 samples.
Framing rests on the assumption that the statistical properties of speech are
approximately constant over short intervals.

\subsection{Energy and Zero-Crossing Rate}

The simplest feature is the \textbf{frame energy}. For frame $t$ with $N$
samples, the log energy is computed as

\begin{equation}
E_t = \log\!\left(\frac{1}{N}\sum_{n=1}^{N} x_t^2[n] + \epsilon\right),
\label{eq:energy}
\end{equation}

where $\epsilon$ is a small positive value that avoids the logarithm of zero.
Energy alone is not enough to detect speech, because loud noise also has high
energy. A complementary feature is the \textbf{zero-crossing rate}, which counts
the sign changes of the signal within the frame:

\begin{equation}
Z_t = \frac{1}{2(N-1)}\sum_{n=2}^{N}
      \bigl|\operatorname{sgn}(x_t[n]) - \operatorname{sgn}(x_t[n-1])\bigr|.
\label{eq:zcr}
\end{equation}

Voiced sounds have a low zero-crossing rate and fricatives a high one; combining
the two features is the basis of classical voice activity
detectors~\cite{sohn1999vad}.

\subsection{Fundamental Frequency and the YIN Method}

The \textbf{fundamental frequency}, $F_0$, is the vibration frequency of the
vocal folds and is perceived as pitch. In this project $F_0$ is estimated with the
YIN method~\cite{decheveigne2002yin}, which rests on the squared difference
function:

\begin{equation}
d_t(\tau) = \sum_{n=1}^{W} \bigl(x_t[n] - x_t[n+\tau]\bigr)^2,
\label{eq:yin-diff}
\end{equation}

where $\tau$ is the lag and $W$ the length of the analysis window. A
``cumulative mean normalized'' version of this function is then built, and the
first local minimum that falls below a threshold is chosen as the period
$\hat{\tau}$; the fundamental frequency is $F_0 = f_s / \hat{\tau}$. If no minimum
falls below the threshold, the frame is marked ``unvoiced''. In this project the
search range of the fundamental frequency is 60--400 Hz, which covers male and
female speakers.

$F_0$ matters for interruption detection because the onset of human speech is
usually accompanied by the appearance of voiced structure, while ambient noise or
loudspeaker playback has no stable pitch pattern.

\subsection{Mel-Frequency Cepstral Coefficients}

MFCCs, the \textbf{mel-frequency cepstral coefficients}~\cite{davis1980mfcc}, are
the most widely used short-term representation of speech; they compress the
frame spectrum in a way that matches human auditory perception. They are computed
as follows:

\begin{enumerate}
\item apply a window (for example a Hamming window) and compute the discrete
      Fourier transform of the frame;
\item compute the power spectrum and pass it through a bank of triangular filters
      whose centres are spread on the mel scale. The mel scale relates to linear
      frequency as
      \begin{equation}
      m(f) = 2595 \log_{10}\!\left(1 + \frac{f}{700}\right);
      \label{eq:mel}
      \end{equation}
\item take the logarithm of each filter output;
\item apply the discrete cosine transform to decorrelate the components and keep
      the first coefficients.
\end{enumerate}

This project keeps 13 coefficients. So that the features also carry information
about \emph{change} over time, the first derivative of these coefficients, known as
the ``delta'', is computed and added to the feature vector. The features of each
frame therefore comprise energy, zero-crossing rate, fundamental frequency, a
voicing flag, 13 MFCCs and the 13 corresponding deltas.

\subsection{Context Window and Decision Vector}

A decision about an interruption cannot rely on a single frame. The features of
the frames of a ``context window'' are therefore aggregated. In this project the
context window is 400 ms, that is, 20 consecutive frames. For each dimension of the
frame feature vector, three statistics (the mean, the standard deviation and the
maximum) are computed over the window and concatenated into a single vector.
This vector is the input of the classifier.

\section{Speech Recognition}
\label{sec:asr}

\textbf{Speech recognition} maps an audio signal to a sequence of words. Current
architectures usually combine a neural encoder with an alignment-free training
objective. Two common objectives are Connectionist Temporal Classification
(CTC)~\cite{graves2006ctc}, which marginalizes over all possible alignments under
a conditional independence assumption, and the RNN Transducer
(RNNT)~\cite{graves2012rnnt}, which also involves a language prediction network in
the decision and is better suited to streaming recognition.

The encoder used in the project's recognition service belongs to the Conformer
family~\cite{gulati2020conformer}, which combines
self-attention~\cite{vaswani2017attention} with convolution: self-attention models
long-range dependencies and convolution local spectral patterns. The model is a
hybrid RNNT-CTC model with subword units, implemented with
NeMo~\cite{kuchaiev2019nemo} and served as a local service.

An alternative approach is large-scale weakly supervised training; its prominent
example is Whisper~\cite{radford2023whisper}, trained on hundreds of thousands of
hours of multilingual data, which this project used only in component comparison
tests.

\section{Text-to-Speech}
\label{sec:tts}

Current \textbf{text-to-speech} systems are usually end-to-end. The
VITS~\cite{kim2021vits} architecture, on which the Piper engine~\cite{piper2023} is
based, combines a conditional variational autoencoder with adversarial training and
a flow-based duration predictor, and generates the waveform directly without a
separate vocoder.

For Persian, text pre-processing matters particularly. Persian script does not
write short vowels, the zero-width non-joiner (ZWNJ) marks the boundaries of the
parts of compound words, and the written form of words can differ across sources.
For this reason the project applies a deterministic Persian text normalizer before
both evaluation and speech generation, so that character forms, ZWNJs and digits
are unified.

Using a \emph{deterministic} speech engine with a fixed voice is a methodological
decision: because the audio output for a given input text is always the same,
measurements based on it become reproducible.

\section{Large Language Models and Low-Rank Adaptation}
\label{sec:llm}

\subsection{Autoregressive Language Models}

A \textbf{language model} factorizes the probability of a token sequence
$y_{1:T}$ as a chain:

\begin{equation}
p_\theta(y_{1:T}) = \prod_{t=1}^{T} p_\theta(y_t \mid y_{<t}),
\label{eq:lm}
\end{equation}

and is fitted by minimizing the cross-entropy loss on training data. The dominant
architecture is the decoder-only Transformer~\cite{vaswani2017attention}.
\emph{Instruction-following} models are built with a supervised fine-tuning stage on
instruction--response pairs~\cite{ouyang2022instructgpt}; the project's text
responder belongs to this class.

\subsection{Teacher Forcing and Exposure Bias}
\label{sec:exposure-bias}

A point that is decisive for interpreting the results of Chapter~\ref{ch:results}
is the difference between training and inference conditions. In training, the
model receives the ``correct'' previous tokens as its condition at every step; this
is called \textbf{teacher forcing}. The loss reported in training and validation is
computed under this condition.

At inference, however, the model must condition on \emph{its own generated}
tokens. If the distribution of generated tokens drifts from that of the reference
tokens, the model enters a region of the state space it has not seen in training,
and errors can grow cumulatively. This mismatch is known as \textbf{exposure
bias}~\cite{ranzato2016scheduledsampling}, and its standard remedy is ``scheduled
sampling'', which gradually replaces reference tokens with generated ones during
training~\cite{bengio2015scheduledsampling}.

The practical consequence of this concept is explicit: ``lower loss under teacher
forcing does not by itself guarantee that the model gives usable output in
autoregressive generation''. The direct-track experiment of this project is a clear
example of this gap.

\subsection{Parameter-Efficient Adaptation and LoRA}

Full fine-tuning of a multi-billion-parameter model on limited hardware is not
practical. \textbf{Parameter-efficient adaptation} methods train only a small
number of parameters~\cite{houlsby2019adapters,li2021prefixtuning}. The method used
in this project is LoRA~\cite{hu2022lora}. Its idea is that the update needed for a
weight matrix has \emph{low rank}. For a frozen weight
$W_0 \in \mathbb{R}^{d \times k}$, instead of learning a full $\Delta W$, two small
matrices $A \in \mathbb{R}^{r \times k}$ and $B \in \mathbb{R}^{d \times r}$ with
$r \ll \min(d,k)$ are trained, and the layer output becomes

\begin{equation}
h = W_0 x + \frac{\alpha}{r} B A x .
\label{eq:lora}
\end{equation}

At the start, $B$ is zero, so the adapted model equals the base model. The
parameter $r$ is the rank and $\alpha$ a scale factor. This project tried ranks 64
and 128. The key advantage of LoRA for this work is that the base weights stay
frozen, so a failed adapter never damages the base model and can simply be
removed.

\section{Neural Audio Codecs and Speech-to-Speech Models}
\label{sec:codec}

\subsection{Residual Vector Quantization}

For a language model to generate speech, audio must be turned into a sequence of
discrete tokens. A \textbf{neural audio codec} does this with an autoencoder with a
Residual Vector Quantizer (RVQ) at its bottleneck. Prominent examples are
SoundStream~\cite{zeghidour2021soundstream} and EnCodec~\cite{defossez2023encodec}.

In residual quantization the latent vector $z$ is approximated successively. If
$Q_i$ is the $i$-th codebook, then with $r_0 = z$:

\begin{equation}
q_i = Q_i(r_{i-1}), \qquad r_i = r_{i-1} - q_i, \qquad
\hat{z} = \sum_{i=1}^{N_q} q_i .
\label{eq:rvq}
\end{equation}

The first codebook carries coarse and mostly ``semantic'' information, and the
following codebooks add ``acoustic'' detail. This hierarchy is why a larger weight
is given to the first codebook's loss in training speech models; this weight was the
changed factor in one of the project's experiments.

\subsection{Language Modelling on Audio Tokens}

Once audio is discrete, the same language modelling machinery can be applied to
audio tokens; this idea is used in AudioLM~\cite{borsos2023audiolm} and
VALL-E~\cite{wang2023valle}. The main challenge is that each time step has not one
token but $N_q$ tokens (one per codebook). The solution is the
RQ-Transformer: a large ``temporal'' Transformer works along the time axis, and a
small ``depth'' Transformer generates the $N_q$ tokens of each time step
autoregressively~\cite{lee2022rqvae}.

\subsection{The Moshi Architecture and the ``Inner Monologue''}
\label{sec:moshi-arch}

Moshi~\cite{defossez2024moshi} models dialogue as speech-to-speech, and three of its
properties matter for this project.

First, the \textbf{Mimi codec}: a streaming audio codec that combines a semantic
representation with an acoustic one in residual codebooks and works with an 80 ms
time step.

Second, \textbf{parallel streams}: the model models two audio streams at the same
time, its own speech and the user's. The explicit notion of a ``turn'' therefore
disappears, and overlap, interruptions and backchannels fit naturally into the
model. This property is what makes Moshi a natural candidate for the full-duplex
requirement.

Third, the \textbf{Inner Monologue}: before generating the audio tokens of each
step, the model predicts time-aligned ``text'' tokens. This intermediate text
noticeably improves the linguistic quality of the generated speech. Its consequence
for debugging is that the text loss and the audio loss can be monitored separately;
Chapter~\ref{ch:results} uses this separation to show where the direct track
failed.

The version of the model used in this project is Moshika: a 7-billion-parameter
model with a female voice, 32 Transformer layers, a latent dimension of 4096 and 16
audio codebooks, eight of which are predicted in the depth Transformer.

\section{Full Duplex and Turn-Taking}
\label{sec:duplex}

In human conversation, turn-taking is a precise and fast process. Classical work on
turn organization showed that the gaps between turns are short and
regular~\cite{sacks1974turntaking}; later engineering reviews report the typical gap
as about two tenths of a second~\cite{skantze2021turntaking}.

From an engineering point of view, three levels of capability should be
distinguished:

\begin{itemize}
\item \textbf{Half duplex}: the system does not listen to the microphone while
      playing its reply. The user has to wait.
\item \textbf{Full duplex with external control}: the microphone is active during
      playback and a \emph{separate} component decides whether to stop playback.
      This is the architecture implemented and evaluated in this project.
\item \textbf{Native full duplex}: the generative model itself models the user's
      stream and its own at the same time, and turn-taking is part of the learned
      distribution, as in Moshi.
\end{itemize}

For the three interaction events labelled in this report, the following
distinction holds: an ``interruption'' is when the second speaker starts talking,
continues and takes the turn; a ``backchannel'' is a short affirming utterance that
does not take the turn; and a ``clean turn'' is when the second speaker starts only
after the first speaker has finished.

\section{Classification with Gradient Boosting}
\label{sec:gbdt}

The project's interruption detector uses a \textbf{gradient-boosted tree
model}~\cite{friedman2001gbm} implemented in
scikit-learn~\cite{pedregosa2011sklearn}. This method builds an ensemble in which
each new tree fits the negative gradient of the loss of the current model:

\begin{equation}
F_m(x) = F_{m-1}(x) + \nu \, h_m(x),
\label{eq:gbm}
\end{equation}

where $h_m$ is the new tree and $\nu$ the learning rate. This model family was
chosen for three reasons: the project definition explicitly asks for a module based
on hand-crafted features; its inference on the CPU fits within the few-millisecond
budget of each audio step; and with a limited number of training events, tree models
are more stable than deep neural networks.

\section{Evaluation Measures and Statistical Concepts}
\label{sec:metrics-prelim}

\subsection{Word and Character Error Rates}

The \textbf{Levenshtein} distance~\cite{levenshtein1966} is the minimum number of
insertions, deletions and substitutions needed to turn one string into another. The
word error rate and character error rate are defined on it. This report gives two
forms of aggregation. The ``micro'' form divides the total number of edits by the
total length of the references:

\begin{equation}
\mathrm{WER}_{\text{micro}} =
\frac{\sum_{i} d\bigl(\hat{w}^{(i)}, w^{(i)}\bigr)}{\sum_{i} \bigl|w^{(i)}\bigr|},
\label{eq:wer}
\end{equation}

and the ``macro'' form is the simple mean of the per-sample error rates. For
Persian, the definition of a word boundary matters: in this project word tokens are
separated by whitespace, and the ZWNJ also counts as a boundary. Character tokens
are Unicode code points, excluding whitespace and ZWNJ. This definition makes the
word error rate in Persian much stricter than the character error rate, a point
interpreted in Chapter~\ref{ch:discussion}.

\subsection{Detector Measures}

For the two-class event ``interruption versus non-interruption'', the confusion
matrix gives four values: true positives, false positives, true negatives and false
negatives. From them the accuracy, the F1 score of the interruption class, the False
Accept Rate (FAR) and the False Reject Rate (FRR) are computed. In full-duplex
systems these two error rates are not equivalent: a false accept interrupts the
system's reply for no reason, while a false reject makes the system ignore the
user's interruption.

\subsection{Confidence Intervals and Group Correlation}

For proportions, the \textbf{Wilson} confidence interval~\cite{wilson1927interval} is
reported; unlike the simple normal approximation, it remains valid for small samples
and for proportions near the boundary.

More importantly, events drawn from one speech session are not \emph{independent};
the speaker, the acoustic conditions and the conversational style are shared. An
event-level confidence interval therefore underestimates the real uncertainty. The
solution used in this project is a \emph{session-level block bootstrap}: in each
resampling repetition, \emph{sessions} are drawn with replacement, not individual
events~\cite{efron1994bootstrap}. Chapter~\ref{ch:results} reports both intervals so
that the effect of this difference can be seen.

\subsection{Group Splits of the Data}
\label{sec:group-split}

If two parts of one speech session fall into the training and the test set, the
model can memorize the speaker or the topic and the test result becomes optimistic.
To avoid this ``group leakage'', the split unit in this project is the source
session identifier: all pairs from one session fall entirely into one and only one
split. The assignment uses deterministic hashing of the session identifier, so that
the split is reproducible.

\subsection{Automatic Judging with a Language Model}
\label{sec:judge}

The semantic quality of a dialogue reply cannot be measured with string-level
measures, because one question has many correct answers. A standard substitute is
to use a large language model as a ``judge''~\cite{zheng2023mtbench}. In this method,
the candidate reply and the comparison replies are given to the judge
\emph{double-blind} and in randomized order, and the judge scores them with a
predefined rubric.

The known limitations of this method are bias towards presentation order, bias
towards reply length, and bias towards replies generated by models of the judge's
own family. This project took three steps to reduce these effects: blinding the
order, deterministic decoding of the judge at temperature zero, and repeating each
call and reporting the agreement rate. Nevertheless, as Chapter~\ref{ch:discussion}
stresses, this measure is an ``automatic approximation'' and no substitute for human
evaluation, especially since in this project the judge and the candidate are of the
same model family.

\subsection{Latency Measures}
\label{sec:latency-prelim}

Two time measures are defined in this report:

\begin{itemize}
\item $T_{\text{first\_audio}}$: the time from the end of the user's speech to the
      first \emph{audible} audio sample of the reply. The official value of this
      measure must be reported by the client, because only the browser knows that
      the audio source was actually scheduled.
\item $T_{\text{barge\_in}}$: the time from the acoustic onset of the user's
      interrupting speech to the moment playback of the reply \emph{actually} stops;
      this value also needs client confirmation.
\end{itemize}

An important methodological point is that ``the server-side time to generate the
whole reply'' is not the $T_{\text{first\_audio}}$ measure. Server generation time is
a useful measurement for debugging, but not the latency acceptance condition.
Chapter~\ref{ch:results} keeps this distinction carefully.
